\section{Loopless Katyusha: a Lyapunov proof in the style of the notes (hypotheses A2, A3)}
\label{sec:lkatyusha}

This note reproduces, in the notation and in the style of Section~5.4.4 of the notes \cite{bach2024learning} (one Lyapunov function, a short list of elementary inequalities, explicit constants), the convergence proof of the \emph{loopless Katyusha} method (L-Katyusha) of Kovalev, Horv\'ath and Richt\'arik \cite{kovalev2020dont}, i.e.\ the accelerated variance-reduced method of Allen-Zhu \cite{allenzhu2018katyusha} in which the outer loop (a full gradient every $m$ iterations) is replaced by a coin flip. The goal is twofold. (A3) We want to see \emph{exactly} which term of the proof is paid by the ``negative momentum'', that is the pull of the extrapolation point towards the reference point $w_t$ at which the full gradient is known. (A2) We want to see why the same proof does not close when Nesterov's momentum of Section~5.2.5 of the notes is combined with the variance-reduced estimator \emph{without} this pull. Every inequality is labelled so that its tightness can be measured in the experiments (Section~\ref{sec:lk-instrument}). The constants are those of \cite{kovalev2020dont} (Algorithm~2, Lemmas~4.1--4.5 and Theorem~4.6 there, proofs in their Appendix~C); we quote them where we use them.

\subsection{Setting, estimator and algorithm in the notation of the notes}
\label{sec:lk-setting}

\begin{assumption}\label{ass:lk}
$F(\theta)=\frac1n\sum_{i=1}^n f_i(\theta)$ on $\R^d$, where each $f_i$ is convex and $R^2$-smooth; we write $L:=R^2$ so that $L$ is the \emph{per-sample} smoothness constant, as in Assumption~1.1 of \cite{kovalev2020dont}. $F$ is $\mu$-strongly convex, with minimizer $\etastar$ and optimal value $F_\ast:=F(\etastar)$. We write $\kappa:=R^2/\mu=L/\mu\ge 1$ and $\sigma:=\mu/L=1/\kappa\in(0,1]$ (the quantity called $\sigma$ in \cite{kovalev2020dont}).
\end{assumption}

Since $\Fp=\frac1n\sum_i \fip$, the function $F$ is itself $L$-smooth (its own constant may be smaller, but we only use $L$). For $\ell_2$-regularized logistic regression, $f_i(\theta)=\log(1+\exp(-y_i x_i^\top\theta))+\frac{\lambda}{2}\norm{\theta}^2$, one can take $L=\max_i\norm{x_i}^2/4+\lambda$ and $\mu=\lambda$ (hypothesis (B3)).

\paragraph{Toolbox.} The proof uses only the following facts, all of which appear in Chapter~5 of the notes. For a convex $L$-smooth $f$ and all $\theta,\eta\in\R^d$:
\begin{itemize}
\item[(T1)] convexity (eq.~(5.7) of the notes): $f(\theta)\ge f(\eta)+\inner{f'(\eta)}{\theta-\eta}$;
\item[(T2)] quadratic upper bound (eq.~(5.10)): $f(\theta)\le f(\eta)+\inner{f'(\eta)}{\theta-\eta}+\frac L2\norm{\theta-\eta}^2$;
\item[(T3)] co-coercivity in its ``Bregman'' form (Prop.~5.4, eq.~(5.12)): $f(\theta)\ge f(\eta)+\inner{f'(\eta)}{\theta-\eta}+\frac1{2L}\norm{f'(\theta)-f'(\eta)}^2$;
\item[(T4)] strong convexity of $F$ (eq.~(5.9)): $F(\theta)\ge F(\eta)+\inner{\Fp(\eta)}{\theta-\eta}+\frac\mu2\norm{\theta-\eta}^2$;
\item[(T5)] three-point identity: $2\inner{a-b}{b-c}=\norm{a-c}^2-\norm{a-b}^2-\norm{b-c}^2$;
\item[(T6)] Young's inequality: for $\beta>0$, $\inner{a}{b}\ge-\frac1{2\beta}\norm a^2-\frac\beta2\norm b^2$;
\item[(T7)] variance is at most the second moment (used as inequality~[I3] in the proof of Prop.~5.9): for a random vector $X$, $\E\norm{X-\E X}^2=\E\norm X^2-\norm{\E X}^2\le\E\norm X^2$.
\end{itemize}
We write $D_F(w,x):=F(w)-F(x)-\inner{\Fp(x)}{w-x}\ge 0$ for the Bregman divergence of $F$; by (T2), $D_F(w,x)\le\frac L2\norm{w-x}^2$.

\begin{definition}[loopless SVRG estimator]\label{def:lsvrg}
Let $p\in(0,1]$. At iteration $t$ the method holds a query point $x_t$ and a \emph{reference point} $w_t$ at which the full gradient $\Fp(w_t)$ has been computed. An index $i(t)$ is drawn uniformly in $\{1,\dots,n\}$ and the estimator is
\begin{equation}\label{eq:lk-estimator}
g_t := f'_{i(t)}(x_t) - f'_{i(t)}(w_t) + \Fp(w_t).
\end{equation}
Independently of $i(t)$ and of the past, a coin $b_t\sim\mathrm{Bernoulli}(p)$ decides whether the reference point is refreshed: $w_{t+1}=y_t$ if $b_t=1$ and $w_{t+1}=w_t$ otherwise (with $y_t$ the ``gradient-step'' sequence defined below). Refreshing costs one pass ($n$ gradients $\fip$) and happens with probability $p$; with $p=1/n$ it costs one gradient per iteration on average.
\end{definition}

\begin{remark}
We write $\Ei$ for the conditional expectation given the state $(y_t,z_t,w_t)$ of the algorithm, i.e.\ over $i(t)$ and $b_t$ (this is the notation $\E[\,\cdot\,|\,y^k,z^k,w^k]$ of \cite{kovalev2020dont}). Since $i(t)$ is uniform, $\Ei[g_t]=\Fp(x_t)-\Fp(w_t)+\Fp(w_t)=\Fp(x_t)$: the estimator is unbiased, exactly as the SAGA direction $\omega_t$ of Prop.~5.9. The difference with SAGA is that the memory is a single point $w_t$ and a single vector $\Fp(w_t)$ instead of the table $(z_i^{(t)})_{i\le n}$, and that each iteration costs two gradients $f'_{i(t)}(x_t),f'_{i(t)}(w_t)$ instead of one.
\end{remark}

\paragraph{The algorithm.} L-Katyusha (Algorithm~2 of \cite{kovalev2020dont}) keeps three sequences $x_t,y_t,z_t$ and the reference point $w_t$. Parameters: $\theta_1,\theta_2>0$ with $\theta_1+\theta_2\le1$, $p\in(0,1]$, $\sigma=\mu/L$ and the step
\begin{equation}\label{eq:lk-eta}
\eta := \frac{\theta_2}{(1+\theta_2)\theta_1}.
\end{equation}
Initialization $y_0=w_0=z_0=\theta_0\in\R^d$ (with $\Fp(w_0)$ computed). For $t=0,1,2,\dots$:
\begin{align}
x_t &= \theta_1 z_t+\theta_2 w_t+(1-\theta_1-\theta_2)y_t, \tag{K1}\label{eq:K1}\\
g_t &= f'_{i(t)}(x_t)-f'_{i(t)}(w_t)+\Fp(w_t), \tag{K2}\label{eq:K2}\\
z_{t+1} &= \frac1{1+\eta\sigma}\Big(\eta\sigma\,x_t+z_t-\frac\eta L\,g_t\Big), \tag{K3}\label{eq:K3}\\
y_{t+1} &= x_t+\theta_1(z_{t+1}-z_t), \tag{K4}\label{eq:K4}\\
w_{t+1} &= \begin{cases} y_t & \text{with probability } p,\\ w_t & \text{with probability } 1-p.\end{cases} \tag{K5}\label{eq:K5}
\end{align}
The reading of the three steps is the one of ``linear coupling'' \cite{allenzhu2017linear}: \eqref{eq:K3} is a \emph{mirror} (or dual-averaging) step with step size $\eta/L$ on a $\mu$-strongly convex surrogate, namely $z_{t+1}=\argmin_z\{\inner{g_t}{z}+\frac L{2\eta}\norm{z-z_t}^2+\frac\mu2\norm{z-x_t}^2\}$ (set the gradient to zero: $g_t+\frac L\eta(z-z_t)+\mu(z-x_t)=0$, multiply by $\eta/L$ and use $\eta\mu/L=\eta\sigma$); \eqref{eq:K4} is a \emph{gradient} step from $x_t$ with step $\theta_1\eta/L$ up to a term of order $\sigma$: indeed by \eqref{eq:K3}, $y_{t+1}=x_t-\frac{\theta_1\eta}{L(1+\eta\sigma)}g_t+\frac{\theta_1\eta\sigma}{1+\eta\sigma}(x_t-z_t)$; and \eqref{eq:K1} is the coupling. The term $\theta_2 w_t$ in \eqref{eq:K1} is the \emph{negative momentum} (``Katyusha momentum'') of \cite{allenzhu2018katyusha}: the query point is pulled towards the point where the estimator has zero variance.

\paragraph{Parameters of the paper.} Theorem~4.6 of \cite{kovalev2020dont} uses
\begin{equation}\label{eq:lk-param}
\theta_2=\tfrac12,\qquad \theta_1=\min\Big\{\sqrt{\tfrac{2\sigma n}{3}},\,\tfrac12\Big\}=\min\Big\{\sqrt{\tfrac{2n}{3\kappa}},\,\tfrac12\Big\},\qquad p=\tfrac1n,\qquad \eta=\frac{\theta_2}{(1+\theta_2)\theta_1}=\frac1{3\theta_1}.
\end{equation}
With this choice the gradient step in \eqref{eq:K4} has length $\theta_1\eta/L=1/(3L)$ (up to the factor $1/(1+\eta\sigma)$), the same as in Katyusha and close to the $1/(4L)$ of SAGA in Prop.~5.9, while the mirror step $\eta/L=1/(3\theta_1L)$ is longer by the factor $1/\theta_1=\max\{\sqrt{3\kappa/(2n)},2\}$: this is where the momentum lives.

\paragraph{The Lyapunov function.} Equation~(13) of \cite{kovalev2020dont} defines, verbatim in their notation, $\Psi^k=\mathcal Z^k+\mathcal Y^k+\mathcal W^k$ with $\mathcal Z^k=\frac{L(1+\eta\sigma)}{2\eta}\|z^k-x^\ast\|^2$, $\mathcal Y^k=\frac1{\theta_1}(f(y^k)-f(x^\ast))$, $\mathcal W^k=\frac{\theta_2(1+\theta_1)}{p\theta_1}(f(w^k)-f(x^\ast))$. In our notation:
\begin{equation}\label{eq:lk-lyap}
\begin{aligned}
\Psi_t &:= Z_t+Y_t+W_t,\\
Z_t&:=\frac{L(1+\eta\sigma)}{2\eta}\norm{z_t-\etastar}^2,\qquad Y_t:=\frac1{\theta_1}\big(F(y_t)-F_\ast\big),\qquad W_t:=\frac{\theta_2(1+\theta_1)}{p\theta_1}\big(F(w_t)-F_\ast\big).
\end{aligned}
\end{equation}
Compare with Prop.~5.9: $V_t=\norm{\theta_t-\etastar}^2+\Delta\sum_i\norm{z_i^{(t)}-\fip(\etastar)}^2$ is a distance plus a memory term; here $Z_t$ is the distance (of the \emph{mirror} sequence), $Y_t$ a function value (of the \emph{gradient} sequence, as in the potential of Exercise~5.14 of the notes for Nesterov's method) and $W_t$ the memory term: it measures how good the reference point is, and it is the loopless substitute for the ``aggregation across the inner loop'' of the original Katyusha proof.

\begin{theorem}[L-Katyusha; Lemma~4.5 and Theorem~4.6 of \cite{kovalev2020dont}]\label{thm:lk}
Under Assumption~\ref{ass:lk}, for any $\theta_1,\theta_2>0$ with $\theta_1+\theta_2\le1$, any $p\in(0,1]$, $\sigma=\mu/L$ and $\eta$ as in \eqref{eq:lk-eta}, the iterates of \eqref{eq:K1}--\eqref{eq:K5} satisfy, for all $t\ge0$,
\begin{equation}\label{ineq:contraction}
\Ei\big[\Psi_{t+1}\big]\;\le\;\frac1{1+\eta\sigma}\,Z_t+\big(1-\theta_1(1-\theta_2)\big)\,Y_t+\Big(1-\frac{p\theta_1}{1+\theta_1}\Big)\,W_t
\qquad\text{(\emph{one-step contraction}).}
\end{equation}
With the parameters \eqref{eq:lk-param}, for all $t\ge0$,
\begin{equation}\label{ineq:rate}
\E\big[\Psi_t\big]\le(1-\rho)^t\,\Psi_0,\qquad \rho:=\min\Big\{\frac{\sigma}{6\theta_1},\frac{\theta_1}{2n}\Big\}
=\begin{cases}\dfrac1{\sqrt{24\,n\kappa}} & \text{if } \kappa\ge\tfrac{8n}{3},\\[2mm] \min\Big\{\dfrac1{3\kappa},\dfrac1{4n}\Big\}\ \ \Big(\ge\dfrac1{8n}\Big) & \text{if } \kappa<\tfrac{8n}{3},\end{cases}
\end{equation}
so that in all cases $\rho\ge\min\{1/(4n),1/\sqrt{24n\kappa}\}$ (\emph{global rate}).
\end{theorem}

The proof occupies Section~\ref{sec:lk-proof}. The complexity in gradient evaluations is derived in Corollary~\ref{cor:lk-complexity}: $\mathcal O\big((n+\sqrt{n\kappa})\log(1/\varepsilon)\big)$, to be compared with $(3n+\frac{16}{3}\kappa)\log(2n/\varepsilon)$ for SAGA (Prop.~5.9) and with $n\sqrt\kappa\log(1/\varepsilon)$ for Nesterov's method of Section~5.2.5 applied to the full gradient: the condition number enters through $\sqrt{n\kappa}$ instead of $\kappa$ (SAGA) or $n\sqrt\kappa$ (accelerated gradient descent).

\begin{remark}[relation with the original Katyusha and with the notes]
Theorem~2.1 of \cite{allenzhu2018katyusha} states (their notation: $\sigma$ is our $\mu$, $m$ the epoch length, and the objective of their Problem~(1.1) is $F=\frac1n\sum_if_i+\psi$ with a proximal term $\psi$) that ``if each $f_i(x)$ is convex, $L$-smooth, and $\psi(x)$ is $\sigma$-strongly convex in Problem~(1.1)'', then Katyusha satisfies $\E[F(\tilde x^S)]-F(x^\ast)\le\mathcal O\big((1+\sqrt{\sigma/(3Lm)})^{-Sm}\big)(F(\tilde x^0)-F(x^\ast))$ when $m\sigma/L\le 3/4$ and $\mathcal O(1.5^{-S})(F(\tilde x^0)-F(x^\ast))$ otherwise, ``choosing $m=\Theta(n)$, Katyusha achieves an $\varepsilon$-additive error using at most $\mathcal O\big((n+\sqrt{nL/\sigma})\log\frac{F(\tilde x^0)-F(x^\ast)}{\varepsilon}\big)$ iterations''. The strong convexity is thus assumed on the proximal term $\psi$ (for $\ell_2$-regularized logistic regression one can take $\psi=\frac\lambda2\norm{\cdot}^2$), whereas L-Katyusha has no proximal term and assumes it on $F$ (Assumption~1.2 of \cite{kovalev2020dont}), as in Assumption~\ref{ass:lk}. The loopless version has the same complexity with a proof that needs a single iteration (no telescoping over an epoch). Allen-Zhu also explains the role of the three sequences: ``In a classical accelerated gradient method, $x^{k+1}$ is only a convex combination of $y^k$ and $z^k$ (or equivalently, $\tau_2=0$ in our formulation)''; we make this precise in Section~\ref{sec:lk-naive} where we show that \eqref{eq:K1}--\eqref{eq:K4} with $\theta_2=0$ and exact gradients is an accelerated gradient method with the rate $(1-\frac12\sqrt{\mu/L})^t$, i.e.\ the method of eq.~(5.16)--(5.17) of the notes up to the choice of constants.
\end{remark}

\subsection{The Lyapunov proof}
\label{sec:lk-proof}

\begin{lemma}[L1: unbiasedness and variance of the estimator; Lemma~4.1 of \cite{kovalev2020dont}]\label{lem:L1}
For all $t$, $\Ei[g_t]=\Fp(x_t)$ and
\begin{align}
\Ei\norm{g_t-\Fp(x_t)}^2 &= \frac1n\sum_{i=1}^n\norm{\fip(x_t)-\fip(w_t)}^2-\norm{\Fp(x_t)-\Fp(w_t)}^2
\;\le\;\frac1n\sum_{i=1}^n\norm{\fip(x_t)-\fip(w_t)}^2 \label{ineq:var-center}\\
&\le\;2L\,\big(F(w_t)-F(x_t)-\inner{\Fp(x_t)}{w_t-x_t}\big)\;=\;2L\,D_F(w_t,x_t). \label{ineq:var-coco}
\end{align}
We call \eqref{ineq:var-center} the \emph{centering inequality} and \eqref{ineq:var-coco} the \emph{co-coercivity bound}; their product is the \emph{variance bound}
\begin{equation}\label{ineq:var}
\Ei\norm{g_t-\Fp(x_t)}^2\le 2L\,D_F(w_t,x_t).
\end{equation}
\end{lemma}

\begin{proof}
Given the state, $x_t$ and $w_t$ are fixed and $X:=f'_{i(t)}(x_t)-f'_{i(t)}(w_t)$ is a random vector with $\Ei X=\Fp(x_t)-\Fp(w_t)$, so that $g_t-\Fp(x_t)=X-\Ei X$ (unbiasedness). By (T7), $\Ei\norm{X-\Ei X}^2=\Ei\norm X^2-\norm{\Ei X}^2$, and $\Ei\norm X^2=\frac1n\sum_i\norm{\fip(x_t)-\fip(w_t)}^2$ because $i(t)$ is uniform: this is \eqref{ineq:var-center}. For \eqref{ineq:var-coco}, apply (T3) to each convex $L$-smooth $f_i$ with $\theta=w_t$, $\eta=x_t$: $\norm{\fip(w_t)-\fip(x_t)}^2\le 2L\big(f_i(w_t)-f_i(x_t)-\inner{\fip(x_t)}{w_t-x_t}\big)$. Average over $i$ and use $\frac1n\sum_i f_i=F$, $\frac1n\sum_i\fip=\Fp$. This is the same use of co-coercivity as in eq.~(5.31) of the notes (there with $\eta=\etastar$, which gives the inner-product form $R^2\,\Fp(\theta)^\top(\theta-\etastar)$). Here we keep the one-sided (Bregman) form (5.12): it is not stronger than the symmetric form (5.11) pointwise, since (5.11) follows from (5.12) by symmetrization, but it is the form that produces the divergence $D_F(w_t,x_t)$, which is exactly what the anchor term cancels in Step~6 of Lemma~\ref{lem:L5}.
\end{proof}

\begin{lemma}[L2: the mirror step; Lemma~4.2 of \cite{kovalev2020dont}]\label{lem:L2}
For every realization of $i(t)$ (no expectation is needed),
\begin{equation}\label{ineq:zstep}
\inner{g_t}{\etastar-z_{t+1}}+\frac\mu2\norm{x_t-\etastar}^2\;\ge\;\frac L{2\eta}\norm{z_t-z_{t+1}}^2+Z_{t+1}-\frac1{1+\eta\sigma}Z_t
\qquad\text{(\emph{mirror-step inequality}),}
\end{equation}
and the slack is exactly $\frac\mu2\norm{x_t-z_{t+1}}^2$ (i.e.\ \eqref{ineq:zstep} becomes an identity if this term is added to the right-hand side).
\end{lemma}

\begin{proof}
Multiplying \eqref{eq:K3} by $1+\eta\sigma$ gives $(1+\eta\sigma)z_{t+1}=\eta\sigma x_t+z_t-\frac\eta L g_t$, i.e.
\begin{equation*}
\frac\eta L\,g_t=\eta\sigma\,(x_t-z_{t+1})+(z_t-z_{t+1}).
\end{equation*}
Take the inner product with $z_{t+1}-\etastar$ and multiply by $L/\eta$, using $L\sigma=\mu$:
\begin{equation*}
\inner{g_t}{z_{t+1}-\etastar}=\mu\inner{x_t-z_{t+1}}{z_{t+1}-\etastar}+\frac L\eta\inner{z_t-z_{t+1}}{z_{t+1}-\etastar}.
\end{equation*}
Apply (T5) twice, with $(a,b,c)=(x_t,z_{t+1},\etastar)$ and $(a,b,c)=(z_t,z_{t+1},\etastar)$:
\begin{align*}
\inner{g_t}{z_{t+1}-\etastar}&=\frac\mu2\Big(\norm{x_t-\etastar}^2-\norm{x_t-z_{t+1}}^2-\norm{z_{t+1}-\etastar}^2\Big)\\
&\quad+\frac L{2\eta}\Big(\norm{z_t-\etastar}^2-\norm{z_t-z_{t+1}}^2-\norm{z_{t+1}-\etastar}^2\Big)\\
&=\frac\mu2\norm{x_t-\etastar}^2-\frac\mu2\norm{x_t-z_{t+1}}^2+\frac L{2\eta}\norm{z_t-\etastar}^2\\
&\quad-\frac L{2\eta}(1+\eta\sigma)\norm{z_{t+1}-\etastar}^2-\frac L{2\eta}\norm{z_t-z_{t+1}}^2,
\end{align*}
where we grouped $\frac\mu2+\frac L{2\eta}=\frac L{2\eta}(\eta\sigma+1)$. By definition \eqref{eq:lk-lyap}, $\frac L{2\eta}(1+\eta\sigma)\norm{z_{t+1}-\etastar}^2=Z_{t+1}$ and $\frac L{2\eta}\norm{z_t-\etastar}^2=\frac{Z_t}{1+\eta\sigma}$. Moving $\inner{g_t}{z_{t+1}-\etastar}$ to the right and $\frac\mu2\norm{x_t-\etastar}^2$ to the left gives the claimed identity with the slack $\frac\mu2\norm{x_t-z_{t+1}}^2\ge0$; dropping it gives \eqref{ineq:zstep}.
\end{proof}

\begin{lemma}[L3: the gradient step; Lemma~4.3 of \cite{kovalev2020dont}]\label{lem:L3}
For every realization of $i(t)$, the \emph{gradient-step inequality} holds:
\begin{equation}\label{ineq:ystep}
\frac1{\theta_1}\Big(F(y_{t+1})-F(x_t)\Big)-\frac{\theta_2}{2L\theta_1}\norm{g_t-\Fp(x_t)}^2\;\le\;\frac L{2\eta}\norm{z_{t+1}-z_t}^2+\inner{g_t}{z_{t+1}-z_t}.
\end{equation}
It is the product of two inequalities, the \emph{smoothness bound}
\begin{equation}\label{ineq:ystep-smooth}
F(y_{t+1})\le F(x_t)+\inner{\Fp(x_t)}{y_{t+1}-x_t}+\frac L2\norm{y_{t+1}-x_t}^2,
\end{equation}
and the \emph{Young bound} on the cross term between the noise of the estimator and the step,
\begin{equation}\label{ineq:ystep-young}
\inner{g_t-\Fp(x_t)}{y_{t+1}-x_t}\;\ge\;-\frac{\theta_2}{2L}\norm{g_t-\Fp(x_t)}^2-\frac L{2\theta_2}\norm{y_{t+1}-x_t}^2.
\end{equation}
\end{lemma}

\begin{proof}
By \eqref{eq:K4}, $y_{t+1}-x_t=\theta_1(z_{t+1}-z_t)$, hence
\begin{equation*}
\frac L{2\eta}\norm{z_{t+1}-z_t}^2+\inner{g_t}{z_{t+1}-z_t}=\frac1{\theta_1}\Big(\frac L{2\eta\theta_1}\norm{y_{t+1}-x_t}^2+\inner{g_t}{y_{t+1}-x_t}\Big).
\end{equation*}
Write $g_t=\Fp(x_t)+(g_t-\Fp(x_t))$ and split the quadratic term as $\frac L{2\eta\theta_1}=\frac L2+\frac L2\big(\frac1{\eta\theta_1}-1\big)$. By the choice \eqref{eq:lk-eta}, $\eta\theta_1=\frac{\theta_2}{1+\theta_2}<1$ and $\frac1{\eta\theta_1}-1=\frac1{\theta_2}$. Therefore the bracket equals
\begin{equation*}
\Big[\frac L2\norm{y_{t+1}-x_t}^2+\inner{\Fp(x_t)}{y_{t+1}-x_t}\Big]+\frac L{2\theta_2}\norm{y_{t+1}-x_t}^2+\inner{g_t-\Fp(x_t)}{y_{t+1}-x_t}.
\end{equation*}
The first bracket is at least $F(y_{t+1})-F(x_t)$ by (T2) applied to $F$ ($\theta=y_{t+1}$, $\eta=x_t$), which is \eqref{ineq:ystep-smooth}. The last two terms are at least $-\frac{\theta_2}{2L}\norm{g_t-\Fp(x_t)}^2$ by (T6) with $\beta=L/\theta_2$, $a=g_t-\Fp(x_t)$, $b=y_{t+1}-x_t$, which is \eqref{ineq:ystep-young}. Dividing by $\theta_1$ gives \eqref{ineq:ystep}.
\end{proof}

\begin{remark}[the step size couples the two momenta]\label{rem:eta-theta2}
In the proof of L3 the coefficient of the variance term is $\frac1{2L}\cdot\frac{\eta\theta_1}{1-\eta\theta_1}$, a function of the step only. The choice \eqref{eq:lk-eta} is precisely the one which makes this coefficient equal to $\frac{\theta_2}{2L}$, with the \emph{same} $\theta_2$ as the weight of $w_t$ in \eqref{eq:K1}. This is what allows the exact cancellation in L5 below. For a general step we write $\nu:=\frac{\eta\theta_1}{1-\eta\theta_1}$ (so $\nu=\theta_2$ for L-Katyusha); L3 then holds with $\theta_2$ replaced by $\nu$ in the variance term. We use this in Section~\ref{sec:lk-naive}.
\end{remark}

\begin{lemma}[L4: the reference point; Lemma~4.4 of \cite{kovalev2020dont}]\label{lem:L4}
For all $t$,
\begin{equation}\label{eq:wstep}
\Ei\big[W_{t+1}\big]=(1-p)\,W_t+\theta_2(1+\theta_1)\,Y_t .
\end{equation}
\end{lemma}

\begin{proof}
By \eqref{eq:K5} and independence of the coin from $i(t)$, $\Ei[F(w_{t+1})]=(1-p)F(w_t)+pF(y_t)$. Subtract $F_\ast$ and multiply by $\frac{\theta_2(1+\theta_1)}{p\theta_1}$: the first term gives $(1-p)W_t$ and the second $\frac{\theta_2(1+\theta_1)}{\theta_1}(F(y_t)-F_\ast)=\theta_2(1+\theta_1)Y_t$.
\end{proof}

\begin{lemma}[L5: combination; Lemma~4.5 of \cite{kovalev2020dont}]\label{lem:L5}
Under the assumptions of Theorem~\ref{thm:lk},
\begin{equation}\label{ineq:zy-step}
\Ei\big[Z_{t+1}+Y_{t+1}\big]\le\frac1{1+\eta\sigma}Z_t+(1-\theta_1-\theta_2)\,Y_t+\frac p{1+\theta_1}\,W_t ,
\end{equation}
and consequently the one-step contraction \eqref{ineq:contraction} holds.
\end{lemma}

\begin{proof}
\emph{Step 1 (strong convexity at the query point).} By (T4) with $\theta=\etastar$ and $\eta=x_t$,
\begin{equation}\label{ineq:sc-x}
F_\ast\ \ge\ F(x_t)+\inner{\Fp(x_t)}{\etastar-x_t}+\frac\mu2\norm{x_t-\etastar}^2
\qquad\text{(\emph{strong convexity at $x_t$}).}
\end{equation}
\emph{Step 2 (decomposition through the coupling).} By \eqref{eq:K1}, $\theta_2w_t+(1-\theta_1-\theta_2)y_t=x_t-\theta_1z_t$, hence
\begin{equation*}
\theta_2(x_t-w_t)+(1-\theta_1-\theta_2)(x_t-y_t)=(1-\theta_1)x_t-(x_t-\theta_1z_t)=\theta_1(z_t-x_t),
\end{equation*}
that is
\begin{equation}\label{eq:xdecomp}
z_t-x_t=\frac{\theta_2}{\theta_1}(x_t-w_t)+\frac{1-\theta_1-\theta_2}{\theta_1}(x_t-y_t).
\end{equation}
Writing $\etastar-x_t=(\etastar-z_t)+(z_t-x_t)$ in \eqref{ineq:sc-x} and using \eqref{eq:xdecomp},
\begin{equation*}
F_\ast\ \ge\ F(x_t)+\frac\mu2\norm{x_t-\etastar}^2+\inner{\Fp(x_t)}{\etastar-z_t}+\frac{\theta_2}{\theta_1}\inner{\Fp(x_t)}{x_t-w_t}+\frac{1-\theta_1-\theta_2}{\theta_1}\inner{\Fp(x_t)}{x_t-y_t}.
\end{equation*}
\emph{Step 3 (convexity towards $y_t$).} Since $1-\theta_1-\theta_2\ge0$, we may use (T1) with $\theta=y_t$, $\eta=x_t$:
\begin{equation}\label{ineq:cvx-y}
\inner{\Fp(x_t)}{x_t-y_t}\ \ge\ F(x_t)-F(y_t)\qquad\text{(\emph{convexity towards $y_t$}).}
\end{equation}
\emph{Step 4 (unbiasedness).} Since $z_t$ is fixed given the state and $\Ei g_t=\Fp(x_t)$,
\begin{equation*}
\inner{\Fp(x_t)}{\etastar-z_t}=\Ei\inner{g_t}{\etastar-z_t}=\Ei\Big[\inner{g_t}{\etastar-z_{t+1}}+\inner{g_t}{z_{t+1}-z_t}\Big].
\end{equation*}
(The deterministic quantity $\frac\mu2\norm{x_t-\etastar}^2$ may be put inside $\Ei$ as well.) Steps~1--4 give
\begin{align*}
F_\ast\ \ge\ & F(x_t)+\frac{\theta_2}{\theta_1}\inner{\Fp(x_t)}{x_t-w_t}+\frac{1-\theta_1-\theta_2}{\theta_1}\big(F(x_t)-F(y_t)\big)\\
&+\Ei\Big[\frac\mu2\norm{x_t-\etastar}^2+\inner{g_t}{\etastar-z_{t+1}}+\inner{g_t}{z_{t+1}-z_t}\Big].
\end{align*}
\emph{Step 5 (the three lemmas).} Inside the expectation, L2 \eqref{ineq:zstep} bounds $\frac\mu2\norm{x_t-\etastar}^2+\inner{g_t}{\etastar-z_{t+1}}$ from below by $\frac L{2\eta}\norm{z_t-z_{t+1}}^2+Z_{t+1}-\frac{Z_t}{1+\eta\sigma}$, then L3 \eqref{ineq:ystep} bounds $\frac L{2\eta}\norm{z_{t+1}-z_t}^2+\inner{g_t}{z_{t+1}-z_t}$ from below by $\frac1{\theta_1}(F(y_{t+1})-F(x_t))-\frac{\theta_2}{2L\theta_1}\norm{g_t-\Fp(x_t)}^2$; finally L1 \eqref{ineq:var} bounds $-\frac{\theta_2}{2L\theta_1}\Ei\norm{g_t-\Fp(x_t)}^2$ from below by $-\frac{\theta_2}{\theta_1}D_F(w_t,x_t)$. Altogether
\begin{align*}
F_\ast\ \ge\ & F(x_t)+\frac{\theta_2}{\theta_1}\inner{\Fp(x_t)}{x_t-w_t}+\frac{1-\theta_1-\theta_2}{\theta_1}\big(F(x_t)-F(y_t)\big)
+\Ei[Z_{t+1}]-\frac{Z_t}{1+\eta\sigma}\\
&+\frac1{\theta_1}\Ei\big[F(y_{t+1})-F(x_t)\big]-\frac{\theta_2}{\theta_1}\Big(F(w_t)-F(x_t)-\inner{\Fp(x_t)}{w_t-x_t}\Big).
\end{align*}
\emph{Step 6 (the cancellation: this is where the negative momentum acts).} The inner products in $w_t$ cancel \emph{exactly}:
\begin{equation}\label{eq:cancel}
\underbrace{\frac{\theta_2}{\theta_1}\inner{\Fp(x_t)}{x_t-w_t}}_{\text{from }\theta_2 w_t\text{ in \eqref{eq:K1}, Step 2}}
\;-\;\underbrace{\frac{\theta_2}{\theta_1}\Big(F(w_t)-F(x_t)-\inner{\Fp(x_t)}{w_t-x_t}\Big)}_{\text{variance term, L1 and L3}}
\;=\;-\frac{\theta_2}{\theta_1}\big(F(w_t)-F(x_t)\big),
\end{equation}
because $-\inner{\Fp(x_t)}{w_t-x_t}=\inner{\Fp(x_t)}{x_t-w_t}$ enters with the sign $-\frac{\theta_2}{\theta_1}$. What survives of the variance term is a \emph{difference of function values} only.\\
\emph{Step 7 (bookkeeping).} Collect the coefficients of $F(x_t)$ after \eqref{eq:cancel}: they sum to $1+\frac{\theta_2}{\theta_1}+\frac{1-\theta_1-\theta_2}{\theta_1}-\frac1{\theta_1}=0$. Hence
\begin{equation*}
F_\ast\ \ge\ -\frac{\theta_2}{\theta_1}F(w_t)-\frac{1-\theta_1-\theta_2}{\theta_1}F(y_t)+\Ei[Z_{t+1}]-\frac{Z_t}{1+\eta\sigma}+\frac1{\theta_1}\Ei[F(y_{t+1})].
\end{equation*}
Since $\frac1{\theta_1}-\frac{\theta_2}{\theta_1}-\frac{1-\theta_1-\theta_2}{\theta_1}=1$, we may write $F_\ast=\frac1{\theta_1}F_\ast-\frac{\theta_2}{\theta_1}F_\ast-\frac{1-\theta_1-\theta_2}{\theta_1}F_\ast$ and rearrange:
\begin{equation*}
\Ei[Z_{t+1}]+\frac1{\theta_1}\Ei\big[F(y_{t+1})-F_\ast\big]\le\frac{Z_t}{1+\eta\sigma}+\frac{1-\theta_1-\theta_2}{\theta_1}\big(F(y_t)-F_\ast\big)+\frac{\theta_2}{\theta_1}\big(F(w_t)-F_\ast\big).
\end{equation*}
By \eqref{eq:lk-lyap}, $\frac{1-\theta_1-\theta_2}{\theta_1}(F(y_t)-F_\ast)=(1-\theta_1-\theta_2)Y_t$ and $\frac{\theta_2}{\theta_1}(F(w_t)-F_\ast)=\frac{p}{1+\theta_1}W_t$: this is \eqref{ineq:zy-step}.\\
\emph{Step 8 (add the reference-point term).} Adding \eqref{eq:wstep} to \eqref{ineq:zy-step}: the coefficient of $Y_t$ becomes $1-\theta_1-\theta_2+\theta_2(1+\theta_1)=1-\theta_1(1-\theta_2)$ and the coefficient of $W_t$ becomes $\frac p{1+\theta_1}+1-p=1-\frac{p\theta_1}{1+\theta_1}$, which is \eqref{ineq:contraction}.
\end{proof}

\begin{proof}[Proof of Theorem~\ref{thm:lk}]
\eqref{ineq:contraction} is Lemma~\ref{lem:L5}. For the rate, \eqref{ineq:contraction} gives $\Ei[\Psi_{t+1}]\le\max\{c_Z,c_Y,c_W\}\Psi_t$ with $c_Z=\frac1{1+\eta\sigma}$, $c_Y=1-\theta_1(1-\theta_2)$, $c_W=1-\frac{p\theta_1}{1+\theta_1}$, and by the tower property $\E[\Psi_{t+1}]\le\max\{c_Z,c_Y,c_W\}\E[\Psi_t]$. With the parameters \eqref{eq:lk-param}:
\begin{itemize}
\item $c_Z\le1-\frac{\sigma}{6\theta_1}$. Indeed $\eta\sigma=\frac{\sigma}{3\theta_1}\le1$: if $\theta_1=\frac12$ then $\eta\sigma=\frac{2\sigma}3\le\frac23$; if $\theta_1=\sqrt{2\sigma n/3}$ then $\eta\sigma=\frac{\sigma}{3\sqrt{2\sigma n/3}}=\sqrt{\frac{\sigma}{6n}}\le1$. And for $x\in[0,1]$, $\frac1{1+x}\le1-\frac x2$ because $(1-\frac x2)(1+x)=1+\frac x2-\frac{x^2}2\ge1$.
\item $c_Y=1-\frac{\theta_1}2$ since $\theta_2=\frac12$.
\item $c_W\le1-\frac{\theta_1}{2n}$: with $p=\frac1n$ and $\theta_1\le\frac12$, $\frac{p\theta_1}{1+\theta_1}=\frac{\theta_1}{n(1+\theta_1)}\ge\frac{2\theta_1}{3n}\ge\frac{\theta_1}{2n}$.
\end{itemize}
Hence $\max\{c_Z,c_Y,c_W\}\le1-\min\{\frac{\sigma}{6\theta_1},\frac{\theta_1}2,\frac{\theta_1}{2n}\}=1-\rho$ with $\rho=\min\{\frac\sigma{6\theta_1},\frac{\theta_1}{2n}\}$, which is the $\theta$ of the proof of Theorem~4.6 of \cite{kovalev2020dont} (``$\theta=\min\{\sigma/6\theta_1,\theta_1/2n\}$''). Explicitly, with $\sigma=1/\kappa$: if $\kappa\ge\frac{8n}3$ then $\frac{2n}{3\kappa}\le\frac14$, $\theta_1=\sqrt{2n/(3\kappa)}$, and $\frac\sigma{6\theta_1}=\frac1{6\kappa}\sqrt{\frac{3\kappa}{2n}}=\frac1{\sqrt{24n\kappa}}$, $\frac{\theta_1}{2n}=\frac1{2n}\sqrt{\frac{2n}{3\kappa}}=\frac1{\sqrt{6n\kappa}}$, so $\rho=\frac1{\sqrt{24n\kappa}}$; if $\kappa<\frac{8n}3$ then $\theta_1=\frac12$ and $\rho=\min\{\frac1{3\kappa},\frac1{4n}\}\ge\frac1{8n}$. Finally $\frac1{3\kappa}\ge\frac1{\sqrt{24n\kappa}}$ iff $9\kappa^2\le24n\kappa$ iff $\kappa\le\frac{8n}3$, which gives the unified bound $\rho\ge\min\{\frac1{4n},\frac1{\sqrt{24n\kappa}}\}$. Iterating, $\E[\Psi_t]\le(1-\rho)^t\Psi_0$.
\end{proof}

\begin{corollary}[gradient complexity]\label{cor:lk-complexity}
With the parameters \eqref{eq:lk-param} and $y_0=z_0=w_0=\theta_0$, for all $t\ge0$,
\begin{equation}\label{ineq:lk-values}
\begin{aligned}
\E\big[F(y_t)-F_\ast\big]&\le\theta_1(1-\rho)^t\Psi_0,\qquad \E\norm{z_t-\etastar}^2\le\frac{2\eta}{L(1+\eta\sigma)}(1-\rho)^t\Psi_0,\\
\theta_1\Psi_0&\le\frac{3L(n+3)}8\norm{\theta_0-\etastar}^2 .
\end{aligned}
\end{equation}
Consequently $\E[F(y_t)-F_\ast]\le\varepsilon$ as soon as $t\ge\big(4n+\sqrt{24n\kappa}\big)\log\frac{3L(n+3)\norm{\theta_0-\etastar}^2}{8\varepsilon}$. An iteration evaluates $f'_{i(t)}(x_t)$ and $f'_{i(t)}(w_t)$ in \eqref{eq:K2} and, with probability $p=1/n$, the $n$ gradients of a full pass: $2+pn=3$ evaluations of some $\fip$ in expectation. If the $n$ gradients $\fip(w)$ computed during the pass are stored ($\mathcal O(nd)$ memory, as for the table of SAGA), $f'_{i(t)}(w_t)$ is a look-up and the count is $1+pn=2$ (this is the ``$\mathcal O(1+pn)$ times per iteration'' of \cite{kovalev2020dont}). Adding the $n$ evaluations of $\Fp(w_0)$ at initialization, the expected number of gradient evaluations is at most
\begin{equation}\label{eq:lk-complexity}
\begin{aligned}
n+\big(12\,n+3\sqrt{24}\,\sqrt{n\kappa}\big)\log\frac{3L(n+3)\norm{\theta_0-\etastar}^2}{8\varepsilon}
&\;\le\;n+\big(12\,n+15\sqrt{n\kappa}\big)\log\frac{3L(n+3)\norm{\theta_0-\etastar}^2}{8\varepsilon}\quad\text{(without storage),}\\
n+\big(8\,n+2\sqrt{24}\,\sqrt{n\kappa}\big)\log\frac{3L(n+3)\norm{\theta_0-\etastar}^2}{8\varepsilon}
&\;\le\;n+\big(8\,n+10\sqrt{n\kappa}\big)\log\frac{3L(n+3)\norm{\theta_0-\etastar}^2}{8\varepsilon}\quad\text{(with storage),}
\end{aligned}
\end{equation}
and in both cases $\mathcal O\big((n+\sqrt{n\kappa})\log\frac{nL\norm{\theta_0-\etastar}^2}{\varepsilon}\big)$.
\end{corollary}

\begin{proof}
The first two bounds follow from $Y_t\le\Psi_t$, $Z_t\le\Psi_t$ and \eqref{ineq:rate}. For $\Psi_0$: $Z_0=\frac{L(1+\eta\sigma)}{2\eta}\norm{\theta_0-\etastar}^2=\frac{3\theta_1L(1+\eta\sigma)}2\norm{\theta_0-\etastar}^2$ since $1/\eta=3\theta_1$; and $Y_0+W_0=\big(\frac1{\theta_1}+\frac{n(1+\theta_1)}{2\theta_1}\big)(F(\theta_0)-F_\ast)\le\frac{2+n(1+\theta_1)}{2\theta_1}\cdot\frac L2\norm{\theta_0-\etastar}^2$ by (T2) at $\eta=\etastar$ ($\Fp(\etastar)=0$). Using $\eta\sigma\le\frac23$, $\theta_1\le\frac12$ and $1\le\frac1{2\theta_1}$: $\Psi_0\le\frac L2\big(\frac52+\frac{2+3n/2}{2\theta_1}\big)\norm{\theta_0-\etastar}^2\le\frac L2\cdot\frac{9/2+3n/2}{2\theta_1}\norm{\theta_0-\etastar}^2=\frac{3L(n+3)}{8\theta_1}\norm{\theta_0-\etastar}^2$. For the iteration count use $(1-\rho)^t\le e^{-\rho t}$ and $1/\rho\le\max\{4n,\sqrt{24n\kappa}\}\le4n+\sqrt{24n\kappa}$.
\end{proof}

\begin{remark}[reading of the rate]
Per iteration, $\Psi$ contracts by $1-\rho$ with $\rho=\min\{\frac\sigma{6\theta_1},\frac{\theta_1}{2n}\}$: the first term is the contraction of the mirror sequence (a strongly convex mirror step with the long step $\eta/L=1/(3\theta_1L)$ contracts the distance by $\frac1{1+\eta\sigma}\approx1-\frac{\sigma}{3\theta_1}$), the second the contraction of the reference term ($w$ is refreshed with probability $p=1/n$ and $Y$ decreases by the fraction $\theta_1(1-\theta_2)=\theta_1/2$). Balancing the two gives $\theta_1\asymp\sqrt{\sigma n}$ and $\rho\asymp\sqrt{\sigma/n}=1/\sqrt{n\kappa}$: this is the $\sqrt{\kappa}$ scaling of hypothesis (A4). For SAGA (Prop.~5.9) the analogous balance is $\min\{\frac1{3n},\frac{3\mu}{16R^2}\}$: the ``parameter'' contraction is $\asymp\sigma$ because the step is $\asymp1/L$; here the mirror step is $\asymp1/(\theta_1L)\gg1/L$ and the price is paid by the variance, which is cancelled by the negative momentum (Step~6) instead of being absorbed by the progress term $-2\gamma(1-\gamma R^2)\Fp(\theta)^\top(\theta-\etastar)$ as in Step~1 of the proof of Prop.~5.9.
\end{remark}

\subsection{Where the negative momentum acts (A3)}
\label{sec:lk-negmom}

The proof of Lemma~\ref{lem:L5} contains a single place where $\theta_2$ is used for something else than bookkeeping: Step~6, equation~\eqref{eq:cancel}. Let us isolate it.

\begin{proposition}[the Bregman-gap cancellation]\label{prop:cancel}
In the chain of Lemma~\ref{lem:L5}, the variance of the estimator enters through L3 and L1 as the term $-\frac{\theta_2}{\theta_1}D_F(w_t,x_t)$ with $D_F(w_t,x_t)=\big(F(w_t)-F(x_t)\big)+\inner{\Fp(x_t)}{x_t-w_t}$, and the weight $\theta_2$ of $w_t$ in \eqref{eq:K1} enters, through the decomposition \eqref{eq:xdecomp}, as the term $+\frac{\theta_2}{\theta_1}\inner{\Fp(x_t)}{x_t-w_t}$. The two inner products cancel identically, and the remainder $-\frac{\theta_2}{\theta_1}(F(w_t)-F(x_t))$ is absorbed as follows: $F(x_t)$ disappears because the coefficients of $F(x_t)$ sum to zero (Step~7), and $\frac{\theta_2}{\theta_1}(F(w_t)-F_\ast)=\frac{p}{1+\theta_1}W_t$ is paid by the reference-point potential, which contracts by \eqref{eq:wstep} because $w$ is refreshed with probability $p$.
\end{proposition}

\begin{proof}
This is Steps~5--8 of the proof of Lemma~\ref{lem:L5}; the identity $D_F(w_t,x_t)=(F(w_t)-F(x_t))+\inner{\Fp(x_t)}{x_t-w_t}$ is the definition of $D_F$.
\end{proof}

Three comments make the mechanism explicit.
\begin{enumerate}
\item \emph{Why an exact cancellation is needed.} The variance term is not small: by (T2), $D_F(w_t,x_t)\le\frac L2\norm{w_t-x_t}^2$ is a squared distance at the scale $L$, whereas the decrease of the potential per step is $\rho\Psi_t$ with $\rho\asymp1/\sqrt{n\kappa}$. To absorb a term $\asymp L\norm{w_t-x_t}^2$ into $\rho\Psi_t$ one would have to convert a squared distance into function values, which by strong convexity costs the factor $\kappa$ (we quantify this in Proposition~\ref{prop:naive-kappa}). The cancellation avoids any such conversion: the inner-product part of $D_F$ is removed identically, and the function-value part (a signed difference of function values; see the discussion of $c(t)$ in Table~\ref{tab:lk-instrument2}) is of the same nature as $Y_t$ and $W_t$.
\item \emph{What the negative momentum costs.} Because $\theta_2w_t$ takes a fraction $\theta_2$ of the weight in \eqref{eq:K1}, the function-value potential $Y_t$ contracts by $\theta_1(1-\theta_2)$ instead of $\theta_1$ (compare \eqref{ineq:contraction} with the deterministic Proposition~\ref{prop:agd-lc} below), and the reference-point potential must carry the weight $\frac{\theta_2(1+\theta_1)}{p\theta_1}\asymp n/\theta_1$ in front of $F(w_t)-F_\ast$. With $\theta_2=\frac12$ this halves the momentum contraction, a constant factor.
\item \emph{Why the step size is tied to $\theta_2$.} By Remark~\ref{rem:eta-theta2}, the coefficient of $\norm{g_t-\Fp(x_t)}^2$ produced by Young's inequality in L3 is $\frac{\nu}{2L\theta_1}$ with $\nu=\frac{\eta\theta_1}{1-\eta\theta_1}$. For the cancellation \eqref{eq:cancel} to be exact one needs $\nu=\theta_2$, i.e.\ \eqref{eq:lk-eta}. A larger step ($\nu>\theta_2$) leaves a residual $+\frac{\nu-\theta_2}{\theta_1}D_F(w_t,x_t)$ on the right-hand side of \eqref{ineq:zy-step}, which is exactly the uncontrolled term of the naive scheme below (with $\nu-\theta_2$ in place of $\nu$). A smaller step ($\nu<\theta_2$) wastes part of the pull.
\end{enumerate}

\subsection{(A2) Naive acceleration: Nesterov's momentum with the SVRG estimator}
\label{sec:lk-naive}

\begin{definition}[naive accelerated SVRG]\label{def:naive}
The \emph{naive scheme} $\mathrm{N}(\theta_1,\eta,p)$ is \eqref{eq:K1}--\eqref{eq:K5} with $\theta_2=0$ and the step $\eta>0$ taken as a free parameter with $\eta\theta_1\le1$:
\begin{equation}\label{eq:naive}
\begin{aligned}
x_t&=\theta_1z_t+(1-\theta_1)y_t,\qquad g_t\text{ as in \eqref{eq:K2}},\qquad z_{t+1}=\frac{\eta\sigma x_t+z_t-\frac\eta Lg_t}{1+\eta\sigma},\\
y_{t+1}&=x_t+\theta_1(z_{t+1}-z_t),\qquad w_{t+1}=\begin{cases}y_t&\text{w.p. }p,\\ w_t&\text{w.p. }1-p.\end{cases}
\end{aligned}
\end{equation}
We set $\nu:=\frac{\eta\theta_1}{1-\eta\theta_1}\in(0,\infty]$, with $\nu=\infty$ in the limiting case $\eta\theta_1=1$; this case is admissible only with exact gradients $g_t=\Fp(x_t)$ (Proposition~\ref{prop:agd-lc}), because the Young step of L3 requires $\nu<\infty$ (Remark~\ref{rem:eta-theta2}). The potential is the same $\Psi_t$ as in \eqref{eq:lk-lyap} (with $W_t$ defined with an arbitrary constant $c_W\ge0$ in place of $\frac{\theta_2(1+\theta_1)}{p\theta_1}$, which is $0$ when $\theta_2=0$).
\end{definition}

\begin{remark}[this is Nesterov's method of the notes with $\Fp$ replaced by $g_t$]\label{rem:naive-book}
Eliminating $z$ from \eqref{eq:naive}: $y_{t+1}=\theta_1z_{t+1}+(1-\theta_1)y_t$ and $x_{t+1}=\theta_1z_{t+1}+(1-\theta_1)y_{t+1}$ give $x_{t+1}=y_{t+1}+(1-\theta_1)(y_{t+1}-y_t)$, which is the extrapolation (5.17) of the notes, $\eta_t=\theta_t+\beta(\theta_t-\theta_{t-1})$, with $\beta=1-\theta_1$. The notes take $\beta=\frac{1-\sqrt{\mu/L}}{1+\sqrt{\mu/L}}=1-2\sqrt{\mu/L}+\mathcal O(\mu/L)$: with $\theta_1=\sqrt{\mu/L}$ the two momentum coefficients agree at zeroth order (both are $1-\Theta(\sqrt{\mu/L})$), but the momentum deficits differ by a factor $2$ ($1-\beta=\sqrt{\mu/L}$ here, against $2\sqrt{\mu/L}+\mathcal O(\mu/L)$ in the notes). The gradient step reads $y_{t+1}=x_t-\frac{\theta_1\eta}{L(1+\eta\sigma)}g_t+\frac{\theta_1\eta\sigma}{1+\eta\sigma}(x_t-z_t)$: it is (5.16) with $\Fp(\eta_{t-1})$ replaced by $g_t$, step $\frac{\theta_1\eta}{L(1+\eta\sigma)}$, and a drift which comes from putting the strong convexity in the mirror step. Since $x_t-z_t=-\frac{1-\theta_1}{\theta_1}(x_t-y_t)$, the drift equals $-\frac{\eta\sigma(1-\theta_1)}{1+\eta\sigma}(x_t-y_t)$, i.e.\ it is of relative size $\eta\sigma$ with respect to the momentum term $x_t-y_t=(1-\theta_1)(y_t-y_{t-1})$, where $\eta\sigma=\sqrt\sigma$ for the parameters of Proposition~\ref{prop:agd-lc} and $\eta\sigma=\frac{\sigma}{3\theta_1}$ for \eqref{eq:lk-param} ($=\sqrt{\sigma/(6n)}$ when $\kappa\ge8n/3$). The drift is therefore small compared with the momentum term, but of the same order $\mathcal O(\sqrt\sigma)$ as the discrepancy between the two momentum coefficients, which is the order that governs the accelerated rate: the two-sequence form of the notes and \eqref{eq:naive} coincide only up to $\mathcal O(\sqrt\sigma)$ corrections of the momentum term. (The drift disappears when the strong convexity is carried by a proximal term instead of the pull towards $x_t$, as in Algorithm~1 of \cite{driggs2022accelerating}, whose $z$-step is $\mathrm{prox}_{\gamma g}(z_t-\gamma g_t)$.) The same scheme with $g_t$ replaced by a mini-batch SVRG estimator is the method of Nitanda \cite{nitanda2014stochastic}; Allen-Zhu writes that it ``corresponds to this paper but \emph{without} Katyusha momentum (i.e., $\tau_2=0$)'' \cite[Appendix~A]{allenzhu2018katyusha}, and Driggs et al.\ write their scheme in the same two-sequence form (``$x_{k+1}=y_k+(1-\tau_k)(y_k-y_{k-1})$'') \cite[Section~2]{driggs2022accelerating}.
\end{remark}

\paragraph{The deterministic case closes.} First we check that the proof machinery of Section~\ref{sec:lk-proof} does prove acceleration for the scheme with $\theta_2=0$ when the gradient is exact; this is the analogue of Exercise~5.14 of the notes in the linear-coupling form of \cite{allenzhu2017linear}.

\begin{proposition}[accelerated gradient descent via L2--L3 with $\theta_2=0$]\label{prop:agd-lc}
Let $g_t=\Fp(x_t)$ in \eqref{eq:naive} (full gradient), $\theta_2=0$, $\eta=1/\theta_1$ (so $\eta\theta_1=1$ and $\nu=\infty$: the limiting case of Definition~\ref{def:naive}) and $\theta_1=\sqrt\sigma=\sqrt{\mu/L}$. Then for all $t$,
\begin{equation}\label{ineq:agd-lc}
Z_{t+1}+Y_{t+1}\le\frac1{1+\sqrt\sigma}Z_t+(1-\sqrt\sigma)\,Y_t\le\Big(1-\frac12\sqrt{\frac\mu L}\Big)(Z_t+Y_t),
\end{equation}
hence $F(y_t)-F_\ast\le\sqrt{\mu/L}\,(1-\frac12\sqrt{\mu/L})^t(Z_0+Y_0)$: the accelerated rate, as for (5.16)--(5.17) in the notes.
\end{proposition}

\begin{proof}
L2 holds verbatim (it does not involve $x_t$'s composition nor the nature of $g_t$). For L3 with $g_t=\Fp(x_t)$ and $\eta\theta_1=1$, the computation of Lemma~\ref{lem:L3} gives $\frac L{2\eta}\norm{z_{t+1}-z_t}^2+\inner{\Fp(x_t)}{z_{t+1}-z_t}=\frac1{\theta_1}\big[\frac L2\norm{y_{t+1}-x_t}^2+\inner{\Fp(x_t)}{y_{t+1}-x_t}\big]\ge\frac1{\theta_1}(F(y_{t+1})-F(x_t))$ by (T2): no Young step and no variance term. Steps~1--5 and 7 of Lemma~\ref{lem:L5} with $\theta_2=0$ (so \eqref{eq:xdecomp} reads $z_t-x_t=\frac{1-\theta_1}{\theta_1}(x_t-y_t)$ and no $w$-terms appear) give $Z_{t+1}+Y_{t+1}\le\frac1{1+\eta\sigma}Z_t+(1-\theta_1)Y_t$ exactly as in \eqref{ineq:zy-step} with $\theta_2=0$, $p=0$. With $\eta=1/\theta_1$ and $\theta_1=\sqrt\sigma$: $\eta\sigma=\sqrt\sigma\le1$, hence $\frac1{1+\sqrt\sigma}\le1-\frac{\sqrt\sigma}2$ and $1-\sqrt\sigma\le1-\frac{\sqrt\sigma}2$.
\end{proof}

\paragraph{The stochastic case does not close.} Now replace $\Fp(x_t)$ by the loopless SVRG estimator.

\begin{proposition}[one-step inequality of the naive scheme]\label{prop:naive}
For the scheme $\mathrm N(\theta_1,\eta,p)$ of Definition~\ref{def:naive} with $\eta\theta_1<1$ (so that $\nu<\infty$), for all $t$,
\begin{equation}\label{ineq:naive}
\begin{aligned}
\Ei\big[Z_{t+1}+Y_{t+1}\big]&\le\frac1{1+\eta\sigma}Z_t+(1-\theta_1)\,Y_t+\frac{\nu}{2L\theta_1}\Ei\norm{g_t-\Fp(x_t)}^2\\
&\le\frac1{1+\eta\sigma}Z_t+(1-\theta_1)\,Y_t+\frac{\nu}{\theta_1}\,D_F(w_t,x_t),
\end{aligned}
\end{equation}
where the last inequality is the variance bound \eqref{ineq:var}. The extra term $\frac\nu{\theta_1}D_F(w_t,x_t)$ is nonnegative; it vanishes only when $F$ is affine on the segment $[x_t,w_t]$.
\end{proposition}

\begin{proof}
L1 and L2 hold verbatim. L3 holds with $\theta_2$ replaced by $\nu$ (Remark~\ref{rem:eta-theta2}; the proof of Lemma~\ref{lem:L3} only used $\frac1{\eta\theta_1}-1=\frac1\nu>0$ in Young's inequality with $\beta=L/\nu$). In the proof of Lemma~\ref{lem:L5}, Step~2 now reads $z_t-x_t=\frac{1-\theta_1}{\theta_1}(x_t-y_t)$, so Steps~1--5 give
\begin{align*}
F_\ast\ \ge\ & F(x_t)+\frac{1-\theta_1}{\theta_1}\big(F(x_t)-F(y_t)\big)+\Ei[Z_{t+1}]-\frac{Z_t}{1+\eta\sigma}\\
&+\frac1{\theta_1}\Ei\big[F(y_{t+1})-F(x_t)\big]-\frac{\nu}{2L\theta_1}\Ei\norm{g_t-\Fp(x_t)}^2 .
\end{align*}
The coefficients of $F(x_t)$ sum to $1+\frac{1-\theta_1}{\theta_1}-\frac1{\theta_1}=0$, and $F_\ast=\frac1{\theta_1}F_\ast-\frac{1-\theta_1}{\theta_1}F_\ast$; rearranging as in Step~7 gives the first inequality of \eqref{ineq:naive}. There is no Step~6: nothing in the chain produces a term $+\frac\nu{\theta_1}\inner{\Fp(x_t)}{x_t-w_t}$, because $w_t$ does not appear in $x_t$. Nonnegativity of $D_F$ is (T1), and $D_F(w_t,x_t)=0$ iff $F(w_t)=F(x_t)+\inner{\Fp(x_t)}{w_t-x_t}$, which for a convex function means that $F$ coincides with its tangent on the segment.
\end{proof}

Thus, compared with \eqref{ineq:zy-step}, the right-hand side of \eqref{ineq:naive} contains the \emph{uncontrolled term} $\frac\nu{\theta_1}D_F(w_t,x_t)$, where $\nu$ is tied to the step: for the Katyusha step ($\eta\theta_1=\frac13$) $\nu=\frac12$, and for the deterministic step of Proposition~\ref{prop:agd-lc} ($\eta\theta_1\to1$) $\nu\to\infty$. The potential $\Psi$ decreases, per step, by at most $\frac{\eta\sigma}{1+\eta\sigma}Z_t$ (distance of the mirror sequence to $\etastar$), $\theta_1Y_t$ (function value at $y_t$) and, if one adds a reference term $c_W(F(w_t)-F_\ast)$, by $pc_W(F(w_t)-F_\ast)-pc_W(F(y_t)-F_\ast)$; none of these involves the query point $x_t$ or $\Fp(x_t)$. The next statement quantifies what it costs to express $D_F(w_t,x_t)$ in terms of the potential.

\begin{proposition}[closing \eqref{ineq:naive} by smoothness and strong convexity costs a factor $\kappa$]\label{prop:naive-kappa}
Under Assumption~\ref{ass:lk}, for the scheme $\mathrm N(\theta_1,\eta,p)$,
\begin{equation}\label{ineq:naive-kappa}
\frac\nu{\theta_1}D_F(w_t,x_t)\;\le\;2\kappa\nu\Big[\frac1{\theta_1}\big(F(w_t)-F_\ast\big)+\frac{\eta}{1+\eta\sigma}Z_t+(1-\theta_1)\,Y_t\Big],
\end{equation}
so that, writing $W_t=c_W(F(w_t)-F_\ast)$ with $c_W>0$ and using \eqref{eq:wstep} with $\theta_2(1+\theta_1)/(p\theta_1)$ replaced by $c_W$,
\begin{equation}\label{ineq:naive-closed}
\Ei[\Psi_{t+1}]\le\Big(\frac{1+2\kappa\nu\eta}{1+\eta\sigma}\Big)Z_t+\Big(1-\theta_1+2\kappa\nu(1-\theta_1)+pc_W\theta_1\Big)Y_t+\Big(1-p+\frac{2\kappa\nu}{\theta_1c_W}\Big)W_t .
\end{equation}
The coefficient of $Z_t$ is $<1$ iff $2\kappa\nu\eta<\eta\sigma$, i.e.\ $\nu<\frac1{2\kappa^2}$; the coefficient of $Y_t$ is $<1$ only if $2\kappa\nu(1-\theta_1)<\theta_1$, i.e.\ $\nu<\frac{\theta_1}{2\kappa(1-\theta_1)}$; the coefficient of $W_t$ is $<1$ iff $c_W>\frac{2\kappa\nu}{p\theta_1}$ (substituting this bound in the $Y$-coefficient sharpens the second condition to $\nu<\frac{\theta_1}{2\kappa(2-\theta_1)}$, which we do not need). Suppose all three coefficients are $<1$. Since $\nu\ge\eta\theta_1$, two cases: if $\theta_1\le\frac12$, the $Y$-condition gives $\eta\le\frac\nu{\theta_1}<\frac1{2\kappa(1-\theta_1)}\le\frac1\kappa$, hence $\eta\sigma<\kappa^{-2}$; if $\theta_1>\frac12$, the $Z$-condition gives $\eta\le\frac\nu{\theta_1}<\frac1{2\kappa^2\theta_1}<\kappa^{-2}$, hence $\eta\sigma<\kappa^{-3}$. In all cases $\eta\sigma<\kappa^{-2}$, and the certified contraction factor is at least the $Z$-coefficient, $\frac{1+2\kappa\nu\eta}{1+\eta\sigma}\ge\frac1{1+\eta\sigma}>\frac1{1+\kappa^{-2}}\ge1-\kappa^{-2}$. Hence no choice of $(\theta_1,\eta,p,c_W)$ makes \eqref{ineq:naive-closed} certify a rate better than $1-\kappa^{-2}$ per iteration: worse than gradient descent with step $1/L$, which contracts by $1-1/\kappa$ per iteration (per gradient evaluation, with $p=1/n$, the comparison is $1-\frac1{3\kappa^2}$ against $1-\frac1{n\kappa}$, so the naive certificate is worse as soon as $\kappa>n/3$), let alone the accelerated $1-\Theta(1/\sqrt{n\kappa})$.
\end{proposition}

\begin{proof}
By (T2), $D_F(w_t,x_t)\le\frac L2\norm{w_t-x_t}^2\le L\norm{w_t-\etastar}^2+L\norm{x_t-\etastar}^2$ (triangle and $(a+b)^2\le2a^2+2b^2$). By (T4) at $\eta=\etastar$, $\norm{\theta-\etastar}^2\le\frac2\mu(F(\theta)-F_\ast)$ for all $\theta$, so $L\norm{w_t-\etastar}^2\le2\kappa(F(w_t)-F_\ast)$ and $L\norm{x_t-\etastar}^2\le2\kappa(F(x_t)-F_\ast)$. By convexity of $F$ and \eqref{eq:naive}, $F(x_t)-F_\ast\le\theta_1(F(z_t)-F_\ast)+(1-\theta_1)(F(y_t)-F_\ast)$, and by (T2) at $\etastar$, $F(z_t)-F_\ast\le\frac L2\norm{z_t-\etastar}^2=\frac{\eta}{1+\eta\sigma}Z_t$ by \eqref{eq:lk-lyap}; also $F(y_t)-F_\ast=\theta_1Y_t$. Multiplying by $\nu/\theta_1$ gives \eqref{ineq:naive-kappa}. Adding \eqref{ineq:naive-kappa} to \eqref{ineq:naive} and then $\Ei[W_{t+1}]=(1-p)W_t+pc_W\theta_1Y_t$ gives \eqref{ineq:naive-closed}; the stated conditions are read off the coefficients, using $\nu\ge\eta\theta_1$ and $\sigma=1/\kappa$.
\end{proof}

\begin{remark}[other routes, and what the statement does and does not say]\label{rem:naive-honest}
(i) The route through $Z_t$ instead of $\mu$ is no better: $L\norm{x_t-\etastar}^2\le L\theta_1\norm{z_t-\etastar}^2+L(1-\theta_1)\norm{y_t-\etastar}^2$ and $L\theta_1\norm{z_t-\etastar}^2=\frac{2\eta\theta_1}{1+\eta\sigma}Z_t$, so the $Z$-coefficient receives $\frac{2\nu\eta}{1+\eta\sigma}$ against the available $\frac{\eta\sigma}{1+\eta\sigma}$: again $\nu<\sigma/2=\frac1{2\kappa}$, a factor $\kappa$ away from the $\nu=\Theta(1)$ needed for a step $\asymp1/L$. (ii) Paying $D_F$ by the reference term alone is impossible: the reference potential pays differences of function values, and $D_F(w_t,x_t)-(F(w_t)-F(x_t))=\inner{\Fp(x_t)}{x_t-w_t}$ is not sign-definite and not a difference of function values; this inner product is precisely what the anchor term produces in \eqref{eq:cancel}. (iii) The comparison with the \emph{non-accelerated} proof of Prop.~5.9 is instructive: there the variance term $2\gamma^2R^2\Fp(\theta_{t-1})^\top(\theta_{t-1}-\etastar)$ (eq.~(5.31)) is absorbed by the progress term $-2\gamma\Fp(\theta_{t-1})^\top(\theta_{t-1}-\etastar)$, which is possible because both scale with the \emph{same} quantity and the step is $\gamma\le1/(4R^2)$; the price is that the contraction is then $\asymp\gamma\mu\asymp\sigma$, not $\sqrt\sigma$. In the accelerated proof the progress is measured by $Y_t$ and $Z_t$, which are not comparable with $D_F(w_t,x_t)$ without the factor $\kappa$. (iv) Propositions~\ref{prop:naive} and~\ref{prop:naive-kappa} show that \emph{this} Lyapunov argument does not certify acceleration for the naive scheme; they do not prove that the scheme is not accelerated (no lower bound specific to the scheme is known to us).
\end{remark}

\paragraph{What is known.} We summarize the literature on the anchor-free scheme, quoting the statements we rely on.
\begin{itemize}
\item \emph{Nitanda (2014, 2016): acceleration needs large mini-batches.} Corollary~1 and Table~1 of \cite{nitanda2014stochastic} give, for Acc-Prox-SVRG with mini-batch size $b$, the overall complexity $\mathcal O\big((n+\frac{n-b}{n-1}\kappa)\log\frac1\varepsilon\big)$ when $b<b_0$ and $\mathcal O\big((n+b\sqrt\kappa)\log\frac1\varepsilon\big)$ otherwise (with a threshold $b_0$ depending on $n$ and $\kappa$), and the paper comments: ``if $b=1$, then Acc-Prox-SVRG has the same complexity as that of Prox-SVRG, while if $b=n$ then the complexity of this method is equal to that of APG''. For $b=1$ this is $\mathcal O((n+\kappa)\log\frac1\varepsilon)$: no acceleration. Allen-Zhu \cite[Remark~5.7]{allenzhu2018katyusha} summarizes the same bound as follows: ``In the simple case $\bar L=L$, Nitanda [46] obtained a total work $\tilde{\mathcal O}\big(n+\frac{n-b}{n-1}\frac L\sigma+b\sqrt{L/\sigma}\big)$ (see their page 7), which also implies a phase transition for $b$. However, this result is no better than ours for all parameters $b\in\{1,2,\dots,n\}$. In fact, in terms of total work, Nitanda's result is even slower than SVRG when $b\le n/2$, and slower than Nesterov's method when $b>n/2$.'' (Here $\sigma$ is our $\mu$.) Allen-Zhu adds (footnote~21): ``Nitanda's method is usually not considered as an accelerated method, since it requires mini-batch size to be very large in order to be accelerated. If mini-batch is large then one can use full-gradient method directly and acceleration is trivial.'' Nitanda's later AMSVRG \cite[Theorem~3]{nitanda2016accelerated} with \emph{increasing} mini-batches has, under the assumptions stated there, an overall complexity ``$\mathcal O\big((n+\frac{n\kappa}{n+\sqrt\kappa})\log\frac1\varepsilon\big)$''; since $\frac{n\kappa}{n+\sqrt\kappa}=\frac12H(\kappa,n\sqrt\kappa)$ with $H$ the harmonic mean, this is $\tilde{\mathcal O}(n+\min\{\kappa,n\sqrt\kappa\})$, the entry of Table~1 there for both Acc-SVRG and AMSVRG in the strongly convex column, while the optimal methods (Acc-SDCA, APCG, SPDC) are listed with $\tilde{\mathcal O}(n+\min\{\kappa,\sqrt{n\kappa}\})$; the paper states that ``for strongly convex problems, AMSVRG is slower than optimal methods''. Neither bound is better than the best of SVRG and accelerated gradient descent.
\item \emph{Allen-Zhu's discussion and experiment.} Footnote~1 of \cite{allenzhu2018katyusha}: ``experimentalists have observed that momentums could sometimes help if stochastic gradient iterations are used. However, the so-obtained methods (1) sometimes fail to converge in an accelerated rate, (2) become unstable and hard to tune, and (3) have no support theory behind them.'' Section~7.1 (``Effectiveness of Katyusha momentum'') compares Katyusha with $\tau_2=\frac12$ and $\tau_2=0$ on ridge regression and concludes: ``the old heuristics (i.e., $\tau_2=0$) sometimes indeed make the method faster after careful parameter tuning. However, for certain tasks such as Figure~1(c), without Katyusha momentum the algorithm does not even enjoy an accelerated convergence rate.''
\item \emph{Hu, Wright and Lessard (2018): the SDP certificate needs the anchor.} In \cite{hu2018dissipativity} the one-iteration inequality of Katyusha is recovered as a dissipation inequality certified by a $6\times6$ linear matrix inequality (their Section~4.3 and Lemma~11). Generalizing the parameters they obtain: ``the one-iteration convergence result (27) holds for any $\frac15\le\tau_2<1$ and $\tau_1\le\min\{\frac{5\tau_2-1}{9\alpha L\tau_2},1-\tau_2\}$''. This is derived for one explicit storage function (their matrix $P$ of eq.~(32)) and fixed multipliers ($\rho^2=\frac1{1+\alpha\sigma}$, $\lambda_1=\alpha/\tau_1$): with this certificate the admissible momentum weight $\tau_1$ shrinks to $0$ as $\tau_2\downarrow\frac15$, and no certificate is given for $\tau_2<\frac15$, in particular not for the anchor-free scheme $\tau_2=0$. The paper does not claim that the LMI is infeasible for $\tau_2<\frac15$ (it stresses that ``Lemma 11 works for general choices of $\zeta$ and $\tau_2$'') and reports no search over other certificates, so the absence of an anchor-free certificate in this family is consistent with, but does not prove, Propositions~\ref{prop:naive}--\ref{prop:naive-kappa}. (The companion jump-system/LMI analysis of SAGA and SVRG without momentum is \cite{hu2017unified}.)
\item \emph{Driggs, Ehrhardt and Sch\"onlieb (2022): anchor-free acceleration with mini-batches.} \cite{driggs2022accelerating} analyze the anchor-free scheme \eqref{eq:naive} up to the placement of the strong convexity, for generic estimators: their Algorithm~1 is $x_{k+1}=\tau z_k+(1-\tau)y_k$, $z_{k+1}=\mathrm{prox}_{\gamma g}(z_k-\gamma\tilde\nabla_{k+1})$, $y_{k+1}=\tau z_{k+1}+(1-\tau)y_k$, with the strong convexity assumed on the proximal term $g$ (``If $g$ is $\mu$-strongly convex''), so that the $z$-step is pulled towards the minimizer of $g$ instead of towards $x_t$ with weight $\eta\sigma$ as in \eqref{eq:K3} (see Remark~\ref{rem:naive-book}). They explain the difficulty in the same terms as above: ``In the analysis of Katyusha, negative momentum cancels part of the variance term, leaving telescoping terms that decrease at an accelerated rate. Without negative momentum, we must handle the variance term differently.'' For the SVRG estimator with mini-batch size $b$ and full gradient with probability $1/p$, their Theorem~14 gives, under the condition $b\le32p^2$ and for $g$ $\mu$-strongly convex, $\E\norm{z_T-x^\ast}^2\le\big(1+\min\{\frac{\sqrt{b\mu}}{4p\sqrt{6L}},\frac b{4p}\}\big)^{-T}K_2$, at an expected cost of ``$n/p+b$'' stochastic gradient evaluations per iteration; with $b=1$ and $p=n$ (admissible since $1\le32n^2$; one full pass every $n$ iterations on average, as in L-Katyusha) this is a contraction $1-\Theta\big(\frac1{n\sqrt{\kappa}}\big)$ per iteration, i.e.\ $\mathcal O(n\sqrt\kappa\log\frac1\varepsilon)$ gradient evaluations: worse than SVRG's $n+\kappa$ unless $\kappa\ge n^2$ and worse than $\sqrt{n\kappa}$ by the factor $\sqrt n$. Their Remark~4 is stated for the SAGA estimator and transferred to SVRG: ``when using the SAGA gradient estimator on a strongly convex objective with $b=\mathcal O(n^{2/3})$ and $\gamma=\mathcal O(1/\sqrt{\mu L})$, Algorithm~1 finds a point satisfying $\E\|z_T-x^\ast\|^2\le\epsilon$ in $\mathcal O(n^{2/3}\sqrt\kappa\log(1/\epsilon))$ iterations. This is compared to the complexity lower bound of $\mathcal O(\sqrt{n\kappa}\log(1/\epsilon))$ achieved by Katyusha [37]. SVRG has a similar complexity when $b=n^{2/3}$ and $p=n^{1/3}$'' (again $b\le32p^2$). So, as for Nitanda, acceleration without the anchor is proved only with mini-batches of size $b\asymp n^{2/3}$ and remains suboptimal by $n^{1/6}$.
\item \emph{Indirect acceleration.} Catalyst \cite{lin2015universal} accelerates SAGA/SVRG as inner solvers with complexity $\tilde{\mathcal O}(\sqrt{n\kappa}\log\frac1\varepsilon)$ when $n\le\kappa$ (their Table~1), with an extra logarithmic factor; this is a different mechanism (proximal-point outer loop), not momentum on the estimator.
\end{itemize}
To our knowledge, no theorem shows that \eqref{eq:naive} with $b=1$ attains $\mathcal O((n+\sqrt{n\kappa})\log\frac1\varepsilon)$, and no lower bound excludes it; the honest statement of (A2) is therefore: \emph{the Lyapunov proof of L-Katyusha closes only because of the cancellation \eqref{eq:cancel}; without the anchor the best available certificates (Nitanda, Driggs et al.) give no improvement over $\min\{n+\kappa,\,n\sqrt\kappa\}$ for $b=1$.}

\paragraph{Empirical test of (A2).} We propose to compare, on the same instances as hypotheses (A4) and (B3) (the hard finite-sum instance with $\kappa\gg n$ of the lower-bound note, and $\ell_2$-regularized logistic regression with $d=40$ Gaussian inputs as in Fig.~5.4 of the notes, $\lambda\in\{R^2/n,R^2/(10n),R^2/(100n)\}$ to vary $\kappa$ at fixed $n$):
\begin{enumerate}
\item SAGA with $\gamma=1/(4R^2)$ (Prop.~5.9) and L-SVRG with $\eta=1/(6L)$, $p=1/n$ (Theorem~3.5 of \cite{kovalev2020dont});
\item L-Katyusha with \eqref{eq:lk-param};
\item the naive scheme \eqref{eq:naive} with the \emph{same} $\theta_1,\eta,p$ as L-Katyusha (so that the only difference is the anchor), with the smaller steps $\eta\in\{\frac1{3\theta_1},\frac1{3},\frac1{3\kappa\theta_1}\}$, and in the two-sequence form of the notes (Remark~\ref{rem:naive-book}) with $\beta=\frac{1-\sqrt{\mu/L}}{1+\sqrt{\mu/L}}$ and step $1/L$;
\item the anchor-free scheme tuned as in Theorem~14 of \cite{driggs2022accelerating} with $b=1$ and $b=\lceil n^{2/3}\rceil$.
\end{enumerate}
For each run we measure the linear rate per gradient evaluation, $\hat\rho:=-\text{slope of }\log(F(y_t)-F_\ast)$ against the number of $\fip$ evaluations, and fit the exponent $\alpha$ in $\hat\rho\propto\kappa^{-\alpha}$ over the three values of $\lambda$. Predictions: (P1) $\alpha\approx1$ for SAGA and L-SVRG, $\alpha\approx\frac12$ for L-Katyusha once $\kappa\ge8n/3$ (regime where $\rho=1/\sqrt{24n\kappa}$), $\alpha\approx0$ for all methods when $\kappa\ll n$ (the $n\log\frac1\varepsilon$ term dominates); (P2) the naive scheme with the Katyusha parameters ($\theta_1=\sqrt{2n/(3\kappa)}$ small, mirror step $\frac1{3\theta_1L}$ large) is non-monotone and may diverge for $\kappa\gg n$, because without the anchor $\norm{x_t-w_t}$, hence the uncontrolled term $\frac\nu{\theta_1}D_F(w_t,x_t)$, is not kept small (this is the instability of Allen-Zhu's footnote~1); (P3) with a step small enough to be stable, $\alpha\approx1$ for the naive scheme (no $\sqrt\kappa$ scaling), and its rate is at best that of L-SVRG; (P4) for the naive scheme the ratio $v(t)$ of the uncontrolled term of \eqref{ineq:naive} to the certified decrease (Section~\ref{sec:lk-instrument}) is frequently $\ge1$, i.e.\ the one-step inequality certifies no decrease at all, while for L-Katyusha the size $q(t)$ of the cancelled inner product relative to the certified decrease $\rho\Psi_t$ is of order one or larger, i.e.\ the term removed by the anchor could not have been absorbed into $\rho\Psi_t$; this is the mechanism behind (P2)--(P3). A positive result for the batch-$n^{2/3}$ variant of item~4 and a negative one for $b=1$ would be consistent with \cite{driggs2022accelerating,nitanda2014stochastic}.

\subsection{What to instrument}
\label{sec:lk-instrument}

Each inequality of the proof is a pair (left-hand side, right-hand side) whose ratio can be computed along a run (with $\etastar$ obtained by a long run of L-Katyusha or SAGA; conditional expectations $\Ei$ are computed exactly by one pass over the data, which is affordable on the instances of (B3)). Tables~\ref{tab:lk-instrument} and~\ref{tab:lk-instrument2} list them with the ratio to record; a ratio close to $1$ means the inequality is tight on the trajectory. The last two rows of Table~\ref{tab:lk-instrument2} concern the naive scheme. Convention: a ratio is not recorded (left undefined) at the steps where its denominator vanishes. This happens systematically at $t=0$ for $r_{\mathrm{cen}},r_{\mathrm{coco}},r_{\mathrm{var}},c,u$ and $q$, because $x_0=w_0=\theta_0$ gives $D_F(w_0,x_0)=0$ and an exact estimator $g_0=\Fp(x_0)$, and it happens for $r_{\mathrm{sc}}$ and $\muEff$ when $x_t=\etastar$. For the Young row of Table~\ref{tab:lk-instrument}, note that the step contains the noise: by \eqref{eq:K3}--\eqref{eq:K4}, $y_{t+1}-x_t=-\alpha\big(g_t-\mu(x_t-z_t)\big)$ with $\alpha:=\frac{\theta_1\eta}{L(1+\eta\sigma)}$ (the drift coefficient $\frac{\theta_1\eta\sigma}{1+\eta\sigma}$ equals $\alpha L\sigma=\alpha\mu$), so that, with $V_t:=\Ei\norm{g_t-\Fp(x_t)}^2$, unbiasedness gives $\Ei[-\inner{g_t-\Fp(x_t)}{y_{t+1}-x_t}]=\alpha V_t$ and $\Ei\norm{y_{t+1}-x_t}^2=\alpha^2\big(V_t+\norm{\Fp(x_t)-\mu(x_t-z_t)}^2\big)$; this gives the closed form of the averaged ratio $\bar r_{\mathrm Y}(t)$ in the table.

\begin{table}[htbp]
\centering\small
\begin{tabular}{p{2.2cm}p{5.6cm}p{7.3cm}}
\hline
Label & Inequality & Tightness ratio to measure (and expected behaviour) \\
\hline
\eqref{ineq:var-center} centering & $\Ei\norm{g_t-\Fp(x_t)}^2\le\frac1n\sum_i\norm{\fip(x_t)-\fip(w_t)}^2$ & $r_{\mathrm{cen}}(t)=1-\dfrac{\norm{\Fp(x_t)-\Fp(w_t)}^2}{\frac1n\sum_i\norm{\fip(x_t)-\fip(w_t)}^2}$; close to $1$ when the per-sample gradient differences are not aligned (expected for logistic regression). \\
\eqref{ineq:var-coco} co-coercivity & $\frac1n\sum_i\norm{\fip(x_t)-\fip(w_t)}^2\le2L\,D_F(w_t,x_t)$ & $r_{\mathrm{coco}}(t)=\dfrac{\frac1n\sum_i\norm{\fip(x_t)-\fip(w_t)}^2}{2L\,D_F(w_t,x_t)}$; for logistic regression expected $\approx L_{\mathrm{eff}}/L\ll1$ where $L$ is the worst-case constant $\max_i\norm{x_i}^2/4+\lambda$ (link with (B2)). \\
\eqref{ineq:var} variance & $\Ei\norm{g_t-\Fp(x_t)}^2\le2L\,D_F(w_t,x_t)$ & $r_{\mathrm{var}}(t)=r_{\mathrm{cen}}(t)\,r_{\mathrm{coco}}(t)$. \\
\eqref{ineq:zstep} mirror step & slack $=\frac\mu2\norm{x_t-z_{t+1}}^2$ & $s_z(t)=\dfrac{\frac\mu2\norm{x_t-z_{t+1}}^2}{\rho\,\Psi_t}$ (slack relative to the certified decrease; small means the dropped term is irrelevant). \\
\eqref{ineq:ystep-smooth} smoothness & $F(y_{t+1})-F(x_t)-\inner{\Fp(x_t)}{y_{t+1}-x_t}\le\frac L2\norm{y_{t+1}-x_t}^2$ & $r_{\mathrm{sm}}(t)=\dfrac{D_F(y_{t+1},x_t)}{\frac L2\norm{y_{t+1}-x_t}^2}\in[0,1]$: local curvature along the step over $L$ (this is the $L_{\mathrm{eff}}$ of (B2)). \\
\eqref{ineq:ystep-young} Young & $\inner{g_t-\Fp(x_t)}{y_{t+1}-x_t}\ge-\frac{\theta_2}{2L}\norm{g_t-\Fp(x_t)}^2-\frac L{2\theta_2}\norm{y_{t+1}-x_t}^2$ & $r_{\mathrm{Y}}(t)=\dfrac{-\inner{g_t-\Fp(x_t)}{y_{t+1}-x_t}}{\frac{\theta_2}{2L}\norm{g_t-\Fp(x_t)}^2+\frac L{2\theta_2}\norm{y_{t+1}-x_t}^2}\in[-1,1]$ per realization, and the ratio of conditional means $\bar r_{\mathrm Y}(t)=\dfrac{\alpha V_t}{\frac{\theta_2}{2L}V_t+\frac L{2\theta_2}\alpha^2\big(V_t+\norm{\Fp(x_t)-\mu(x_t-z_t)}^2\big)}$, with $V_t$ and $\alpha$ as in the text above. With \eqref{eq:lk-param}, $\bar r_{\mathrm Y}(t)\le\frac{12(1+\eta\sigma)}{9(1+\eta\sigma)^2+4}\le\frac{12}{13}$, with near-equality in the noise-dominated regime $V_t\gg\norm{\Fp(x_t)-\mu(x_t-z_t)}^2$: the Young step is then nearly tight, and it is loose only when the full gradient or the drift dominates the step. \\
\eqref{ineq:sc-x} strong convexity & $F_\ast-F(x_t)-\inner{\Fp(x_t)}{\etastar-x_t}\ge\frac\mu2\norm{x_t-\etastar}^2$ & $r_{\mathrm{sc}}(t)=\dfrac{2D_F(\etastar,x_t)}{\mu\norm{x_t-\etastar}^2}\ge1$: the Bregman analogue of the $\muEff/\mu$ of (B2). For comparability with (B2) also record the inner-product form $\muEff(x_t)/\mu$ with $\muEff(x_t):=\inner{\Fp(x_t)}{x_t-\etastar}/\norm{x_t-\etastar}^2\ge\mu$; the two differ, since $D_F(\etastar,x_t)\le\inner{\Fp(x_t)}{x_t-\etastar}$ (because $F_\ast\le F(x_t)$), so that $r_{\mathrm{sc}}(t)\le2\muEff(x_t)/\mu$. \\
\eqref{ineq:cvx-y} convexity & $\inner{\Fp(x_t)}{x_t-y_t}\ge F(x_t)-F(y_t)$ & $s_{\mathrm{cvx}}(t)=\dfrac{D_F(y_t,x_t)}{F(y_t)-F_\ast}$ (slack relative to the $Y$-potential). \\
\hline
\end{tabular}
\caption{Inequalities of the one-step lemmas L1--L3 and of the convexity steps of L5, with the ratio to record along a trajectory.}
\label{tab:lk-instrument}
\end{table}

\begin{table}[htbp]
\centering\small
\begin{tabular}{p{2.2cm}p{5.6cm}p{7.3cm}}
\hline
Label & Inequality & Tightness ratio to measure (and expected behaviour) \\
\hline
\eqref{eq:cancel} cancellation & identity & $c(t)=\dfrac{\inner{\Fp(x_t)}{x_t-w_t}}{D_F(w_t,x_t)}\in\R$ and $1-c(t)=\dfrac{F(w_t)-F(x_t)}{D_F(w_t,x_t)}$: the \emph{signed} shares of the Bregman gap removed by the anchor and left as a difference of function values. They sum to $1$ by definition of $D_F$ but are not fractions in $[0,1]$: $c(t)<0$ when $\inner{\Fp(x_t)}{x_t-w_t}<0$ (the anchor term then adds to the variance budget on that step), and $c(t)>1$ when $F(w_t)<F(x_t)$ (the reference point is better than the query point and the function-value share is negative). Record $c(t)$ together with the sign of $F(w_t)-F(x_t)$, and the size of the cancelled term relative to the certified decrease, $q(t)=\dfrac{\big|\frac{\theta_2}{\theta_1}\inner{\Fp(x_t)}{x_t-w_t}\big|}{\rho\,\Psi_t}$ ($q(t)\gtrsim1$ means that the cancelled term could not have been absorbed into the decrease); and $u(t)=\dfrac{\frac{\theta_2}{2L\theta_1}\Ei\norm{g_t-\Fp(x_t)}^2}{\frac{\theta_2}{\theta_1}D_F(w_t,x_t)}=r_{\mathrm{var}}(t)$, the fraction of the cancelled budget actually used by the noise. \\
\eqref{ineq:contraction} contraction & $\Ei\Psi_{t+1}\le\max\{c_Z,c_Y,c_W\}\Psi_t$ & $r_\Psi(t)=\dfrac{\Ei\Psi_{t+1}}{(1-\rho)\Psi_t}$ and the three component ratios $\Ei Z_{t+1}/Z_t$, $\Ei Y_{t+1}/Y_t$, $\Ei W_{t+1}/W_t$ against $c_Z,c_Y,c_W$ (which of $Z,Y,W$ is the bottleneck). \\
\eqref{ineq:rate} global rate & $\E\Psi_t\le(1-\rho)^t\Psi_0$ & $\hat\rho_T(t)=-\frac1T\log\frac{\Psi_{t+T}}{\Psi_t}$ against $\rho$ (expected $\hat\rho\gg\rho$ on logistic regression, cf.\ $r_{\mathrm{coco}},r_{\mathrm{sc}}$). \\
\eqref{ineq:naive} naive, uncontrolled term & $\Ei[Z_{t+1}+Y_{t+1}]\le\frac{Z_t}{1+\eta\sigma}+(1-\theta_1)Y_t+\frac\nu{\theta_1}D_F(w_t,x_t)$ & $v(t)=\dfrac{\frac\nu{\theta_1}D_F(w_t,x_t)}{\frac{\eta\sigma}{1+\eta\sigma}Z_t+\theta_1Y_t}$, the uncontrolled term over the certified decrease of $Z+Y$: $v(t)\ge1$ exactly when the one-step inequality certifies no decrease of $Z_t+Y_t$ at step $t$. \\
\eqref{ineq:naive-kappa} naive, $\kappa$-cost & $\frac\nu{\theta_1}D_F\le2\kappa\nu[\dots]$ & ratio of the two sides; expected $\ll1$ (the bound is loose by $\approx\kappa$), confirming that the obstruction is the proof, not the quantity. \\
\hline
\end{tabular}
\caption{The cancellation, the contraction of the potential, the global rate, and the two inequalities of the naive scheme, with the ratio to record along a trajectory.}
\label{tab:lk-instrument2}
\end{table}

\subsection{Complexity comparison}
\label{sec:lk-complexity-table}

Table~\ref{tab:lk-complexity} collects the number of evaluations of a single $\fip$ needed to reach accuracy $\varepsilon$ (on $\E[F-F_\ast]$ or $\E\norm{\cdot-\etastar}^2$, which differ by factors $L$ or $\mu$ inside the logarithm), with the constants of the cited statements; $\kappa=L/\mu$ with $L$ the per-sample smoothness constant, as everywhere in this note. Counting convention: the evaluation of $\fip$ at a stored reference point ($w_t$ for L-SVRG and L-Katyusha, $\tilde x$ for Katyusha) counts as one gradient unless the $n$ reference gradients are stored when the full gradient is computed ($\mathcal O(nd)$ memory, as for the table of SAGA), in which case it is a look-up; where the two conventions differ we give both counts, the first one without storage.

\begin{table}[htbp]
\centering\small
\begin{tabular}{p{3.4cm}p{6.4cm}p{5.3cm}}
\hline
Method & Gradient evaluations to accuracy $\varepsilon$ & Source and constants \\
\hline
Gradient descent, $\gamma=1/R^2$ & $n\,\kappa\log\frac1\varepsilon$ & Notes, Section~5.4.4 (``running time $dn\frac{R^2}\mu\log\frac1\varepsilon$''); with $\gamma=1/L_F$, $n\frac{L_F}\mu\log\frac1\varepsilon$ where $L_F\le R^2$. \\
Accelerated GD (5.16)--(5.17) & $n\sqrt{\kappa}\log\frac1\varepsilon$ & Notes, Section~5.2.5: rate $(1-\sqrt{\mu/L})^t$; Proposition~\ref{prop:agd-lc}: $(1-\frac12\sqrt{\mu/L})^t$. \\
SGD, $\gamma_t=1/(\mu t)$ & $\mathcal O(\kappa/\varepsilon)$ (no logarithm, sublinear) & Notes, Section~5.4.4 (``the convergence rate has been shown to be $\mathcal O(\kappa/t)$ in Section~5.4.1'') and the table of Section~5.5 (smooth case, rate $L/(t\mu)$); Prop.~5.8 is the Lipschitz case, $2B^2(1+\log t)/(\mu t)$. \\
SAGA, $\gamma=1/(4R^2)$ & $\big(3n+\frac{16}3\kappa\big)\log\frac{2n}\varepsilon$ & Prop.~5.9 of the notes: contraction $1-\min\{\frac1{3n},\frac{3\mu}{16R^2}\}$, one gradient per iteration; \cite{defazio2014saga}. \\
SVRG (epoch length $m$) & $\mathcal O\big((n+\kappa)\log\frac1\varepsilon\big)$ with $m\asymp\kappa$ & \cite{johnson2013accelerating}; as quoted in \cite[Section~1]{kovalev2020dont}: ``the iteration complexity of SVRG is $\mathcal O((n+L/\mu)\log1/\varepsilon)$''. \\
L-SVRG, $\eta=1/(6L)$, $p=1/n$ & $3\max\{6\kappa,2n\}\log\frac{\Phi_0}\varepsilon$ without storage, $2\max\{6\kappa,2n\}\log\frac{\Phi_0}\varepsilon$ with & \cite[Theorem~3.5]{kovalev2020dont}: $\E\Phi^k\le\max\{1-\frac\mu{6L},1-\frac1{2n}\}^k\Phi^0$, $\Phi^k=\|x^k-x^\ast\|^2+\frac{4\eta^2}{pn}\sum_i\|\nabla f_i(w^k)-\nabla f_i(x^\ast)\|^2$; $2+pn=3$ gradients per iteration without storage, $1+pn=2$ with. \\
Katyusha, $m=2n$ & $\mathcal O\big((n+\sqrt{n\kappa})\log\frac{F(\theta_0)-F_\ast}\varepsilon\big)$ & \cite[Theorem~2.1, Remark~2.2 and footnote~11]{allenzhu2018katyusha}: $2.5$ stochastic gradients per iteration without storage, $1.5$ with (quoted below). \\
L-Katyusha, \eqref{eq:lk-param} & $n+\big(12n+15\sqrt{n\kappa}\big)\log\frac{3L(n+3)\norm{\theta_0-\etastar}^2}{8\varepsilon}$ without storage, $n+\big(8n+10\sqrt{n\kappa}\big)\log\frac{3L(n+3)\norm{\theta_0-\etastar}^2}{8\varepsilon}$ with & Corollary~\ref{cor:lk-complexity}; \cite[Theorem~4.6]{kovalev2020dont}: $\rho=\min\{\frac\sigma{6\theta_1},\frac{\theta_1}{2n}\}$; $2+pn=3$ gradients per iteration without storage, $1+pn=2$ with (quoted below); the leading $n$ is the initial $\Fp(w_0)$. \\
Catalyst on SAGA/SVRG & $\tilde{\mathcal O}\big(\sqrt{n\kappa}\log\frac1\varepsilon\big)$ for $n\le\kappa$ (extra $\log\kappa$ factor) & \cite[Table~1]{lin2015universal}. \\
Naive AGD+SVRG, $b=1$ & no bound better than $\min\{n+\kappa,\,n\sqrt\kappa\}\log\frac1\varepsilon$: Nitanda $\tilde{\mathcal O}(n+\frac{n-b}{n-1}\kappa+b\sqrt\kappa)$; Driggs et al.\ $\mathcal O(n\sqrt\kappa\log\frac1\varepsilon)$ & \cite[Remark~5.7]{allenzhu2018katyusha} on \cite{nitanda2014stochastic}; \cite[Theorem~14]{driggs2022accelerating} with $b=1$, $p=n$ (Section~\ref{sec:lk-naive}). \\
Lower bound, randomized incremental gradient methods (any sampling distribution) & $\Omega\big((n+\sqrt{n\kappa})\log\frac{\norm{\theta_0-\etastar}^2}\varepsilon\big)$, for a large enough dimension & \cite[Corollary~3]{lan2018optimal} (quoted below; their $m$ is our $n$, and in their sum convention $L=\sum_iL_i$ is our $L$ after rescaling $f_i\to f_i/n$). \\
Lower bound, any randomized algorithm, $\kappa>161\,n$ & $\Omega\big(n+\sqrt{n\kappa}\log\big(\frac{\varepsilon_0}{\varepsilon}\sqrt{\frac n\kappa}\big)\big)$, with no $n\log\frac1\varepsilon$ term & \cite[Theorem~8]{woodworth2016tight} (quoted below); see hypothesis (A1) and \cite{hannah2018breaking} for the $n\log\frac1\varepsilon$ term and the role of the span assumption. \\
\hline
\end{tabular}
\caption{Gradient complexity of the methods discussed in the notes and in this note, in terms of $n$, $\kappa=L/\mu$ and $\varepsilon$.}
\label{tab:lk-complexity}
\end{table}

\paragraph{Quotations used in Table~\ref{tab:lk-complexity}.} Gradient counts. \cite[footnote~11]{allenzhu2018katyusha}: ``The claim `SVRG or Katyusha computes 1.5 stochastic gradients' requires one to store $\nabla f_i(\tilde x)$ in the memory for each $i\in[n]$, and this costs $O(dn)$ space in the most general setting. If one does not store $\nabla f_i(\tilde x)$ in the memory, then each iteration of SVRG or Katyusha computes 2.5 stochastic gradients for the choice $m=2n$.'' \cite[Section~4]{kovalev2020dont}: ``Algorithm 2 calls stochastic gradient oracle $\mathcal O(1+pn)$ times per iteration in expectation.'' Lower bounds. \cite[Corollary~3]{lan2018optimal}: ``The number of gradient evaluations performed by any randomized incremental gradient methods for finding a solution $\bar x\in X$ of problem (1.1) such that $\E[\|\bar x-x^\ast\|_2^2]\le\epsilon$ cannot be smaller than $\Omega\big((\sqrt{mC}+m)\log\frac{\|x^0-x^\ast\|_2^2}{\epsilon}\big)$ if $n$ is sufficiently large, where $C=L/\mu$ and $L=\sum_{i=1}^mL_i$'' (their $n$ is the dimension), the worst-case instance being analyzed ``for any arbitrary distribution used to choose $f_i$ at each iteration''. \cite[Theorem~8]{woodworth2016tight}, for $m$ components that are $\gamma$-smooth and $\lambda$-strongly convex with $\gamma/\lambda>161m$ and initial gap $\epsilon_0$: any randomized algorithm $\mathcal A$ ``must make $\Omega\big(m+\sqrt{m\frac\gamma\lambda}\log\big(\frac{\epsilon_0}{\epsilon}\sqrt{\frac{m\lambda}{\gamma}}\big)\big)$ queries to $h_F$''.

In the regime $\kappa\gg n$ of hypothesis (A4), the rows read $n\kappa$ (GD), $n\sqrt\kappa$ (AGD), $\kappa$ (SAGA, L-SVRG), $\sqrt{n\kappa}$ (L-Katyusha, optimal): among the incremental methods ($\mathcal O(1)$ gradients per iteration) the exponent of $\kappa$ drops from $1$ to $\frac12$ only for the one whose proof contains the cancellation \eqref{eq:cancel}; accelerated gradient descent also has the exponent $\frac12$ but pays the factor $n$ per iteration.
